% Generated draft: requires mathematical and source review.
\documentclass[11pt]{article}
\usepackage[margin=25mm]{geometry}
\usepackage{fontspec}
\setmainfont{Latin Modern Roman}
\usepackage{amsmath,amssymb,mathtools,unicode-math}
\setmathfont{Latin Modern Math}
\usepackage{enumitem,needspace,xurl,hyperref}
\setlength{\parindent}{0pt}
\setlength{\parskip}{5pt}
\setlength{\emergencystretch}{3em}
\newcommand{\N}{\mathbb{N}}
\newcommand{\Z}{\mathbb{Z}}
\newcommand{\Q}{\mathbb{Q}}
\newcommand{\R}{\mathbb{R}}
\newcommand{\C}{\mathbb{C}}
\newcommand{\F}{\mathbb{F}}
\newcommand{\A}{\mathbb{A}}
\newcommand{\PP}{\mathbb{P}}
\newcommand{\E}{\mathbb{E}}
\newcommand{\eps}{\varepsilon}
\newcommand{\dd}{\,\mathrm{d}}
\begin{document}
\section*{UnsolvedMath ID: 10400080}
\textbf{Source problem code:} catalogue entry
\subsection{Definitions and mathematical statement}
In the standard Kauffman-bracket convention, let $R=\Z[A,A^{-1}]$ and let $S_{2,\infty}(M)$ be the $R$-module generated by isotopy classes of unoriented framed links in a three-manifold $M$, including the empty link, modulo the local skein relation
\[
[L_\times]=A[L_0]+A^{-1}[L_\infty]
\]
and the trivial-circle relation $[L\sqcup\bigcirc]=-(A^2+A^{-2})[L]$, with compatible local pictures/framing conventions. An element $x\ne0$ is torsion if $f(A)x=0$ for some nonzero $f(A)\in R$.

The source, attributed to J. Przytycki, reports that incompressible tori and two-spheres in $M$ yield torsion and asks whether other surfaces can yield torsion as well. The exact hypotheses and construction behind that reported fact are deferred to reference [339], which is absent. In particular essentiality and separation/embedding conditions cannot be inferred for arbitrary spheres or tori from the short sentence.

The requested extension concerns a geometric source of torsion, not merely a manifold that both contains another surface and has torsion from an unrelated torus or sphere. The source does not specify which additional surface types, their incompressibility or proper/closed status, or what construction constitutes yielding torsion. A precise positive instance would identify the surface-induced relation, its nonzero annihilating polynomial and the nonvanishing of the module class. No universal assertion that every higher-genus surface creates torsion is supplied.

The coefficient ring and manifold category are not reproduced. The displayed Laurent-polynomial skein convention is explicit standard context; specialization or changing coefficients can change the torsion question and is not silently allowed.
\subsection{Short English statement}
Can surfaces beyond the source's torus and sphere cases geometrically produce torsion in the Kauffman-bracket skein module? Specify the surface conditions and exhibit a nonzero annihilated class, rather than only coexisting torsion.
\subsection{Sources}
\begin{itemize}
\item Local UnsolvedMath registry, permanent ID 10400080, exact source code catalogue entry, Problem 4.2; source attribution J. Przytycki, no proposer/year in metadata. Registry status open, classification OPEN-TRIAGE; title truncated.
\item Ohtsuki, \emph{Problems on invariants of knots and 3-manifolds} (2002), Problem 4.2, PDF page 74, \url{https://msp.org/gtm/2002/04/gtm-2002-04-024s.pdf}. Supplied pdf-font-structure extraction accessed July 29, 2026; provenance distinguishes 231 bold numbered items. No PDF was accessed here.
\item The source cites [339] for the torus/sphere torsion statement, but its bibliography entry and precise theorem hypotheses are absent. The supplied summary gives no later construction or resolution.
\item Provenance flags current primary-literature status and redistribution terms NEEDS\_REVIEW.
\end{itemize}
\end{document}
